\section{Bab 6\\Penutup}
\addcontentsline{toc}{section}{Bab 6. Penutup}

\subsection{Kesimpulan}
\label{sec:6-1}

Bantar Gebang yang \textit{overload} (rata-rata \textpm{}7.354 ton/hari pada 2025, puncak \textpm{}8.000 ton/hari) dan RDF Rorotan yang berkapasitas 2.200--2.400 ton/hari namun baru satu lini beroperasi karena kendala logistik \parencite{dlhjakarta2026,kompasrorotan2026} menunjukkan bahwa mata rantai yang hilang adalah sistem yang mengarahkan sampah bermutu kalor ke pabrik secara konsisten. \ProductName{} menjawabnya sebagai \textit{execution and verification layer} B2B/B2G SaaS \textit{asset-light} yang mengintegrasikan sensor IoT \textit{fill-level}, \textit{Dynamic Routing AI}, dan \textit{Predictive RDF Allocation System} --- dengan kejujuran teknis bahwa perangkat mengukur kepenuhan dan sistem mengestimasi kelayakan RDF sebagai \textit{proxy classifier} yang dikalibrasi pada fase pilot --- selaras dengan Ingub 5/2026 dan Pergub 102/2021 \parencite{ingub5_2026,pergub102}. Proyeksi keuangan menunjukkan kelayakan skenario dasar (modal awal \textpm{}Rp528,5 juta; NPV \textpm{}Rp975 juta pada diskonto 12\%; IRR \textpm{}31,6\%; NPV tetap positif hingga diskonto 30\%), namun pada skenario konservatif (pendapatan $-$20\%, biaya $+$20\%) proyek belum layak, sehingga pencapaian target akuisisi pada Bab 1.2 dan 3.2 menjadi kunci keberhasilan. Prioritas tim selanjutnya adalah mengajukan percontohan bersama DLH DKI Jakarta dan mengakuisisi lima mitra B2B pertama pada tahun pertama, sebagai dasar pembuktian model sebelum ekspansi ke skala kota penuh.
